% !TeX root = ../guide
\chapter{Dia: Structured Diagramming and Flowcharts}\label{ch:dia}
\chaptermark{Dia}
	\section{Introduction}\label{sec:DiaIntroduction}
		A technical thesis often requires abstract visual representations of your work. 
		Whether you need to map out the logic for a PLC program, draft a block diagram for a hydraulic circuit, or outline the step-by-step workflow of a process, you need a dedicated diagramming tool.

		\textbf{Dia} is a free, open-source drawing program designed specifically for structured diagrams. 
		It serves as an excellent, lightweight alternative to commercial software like Microsoft Visio. 
		This chapter will introduce Dia and explain how to create professional, scalable diagrams that integrate perfectly into your \LaTeX\ document.

	\section{Key Features of Dia}
		Dia is not a freehand drawing tool; it is built for precision, structure, and connectivity.

		\subsection{Extensive Shape Libraries}
			Dia comes pre-loaded with standardized shape libraries for various technical disciplines. 
			You will find dedicated sheets for basic flowcharts, UML (Unified Modeling Language) for software engineering, electrical circuits, pneumatic systems, and network topologies. 
			This ensures your diagrams use industry-standard iconography.

		\subsection{Dynamic Routing and Connectivity}
			When drafting a complex control logic diagram, the last thing you want is for your lines to break when you move a box. 
			Dia features dynamic connection points. 
			When you anchor a line between two blocks and subsequently reposition one of the blocks, the connecting line automatically stretches and reroutes itself, keeping your system architecture intact.

		\subsection{Snap-to-Grid Alignment}
			Professional diagrams require perfect alignment. 
			Dia includes a robust grid system and alignment tools that ensure your blocks and flowcharts are perfectly symmetrical, which is critical for the visual appeal of your printed thesis.

	\section{Getting Started with Dia}
		Dia is incredibly lightweight and easy to get running on any system.

		\subsection{Installation}
			Dia is available for Windows, macOS, and Linux.
			\begin{enumerate}
				\item For Windows users, the easiest way to download the software is via the installer website: \url{http://dia-installer.de/}.
				\item Download the latest installer package.
				\item Follow the standard installation prompts.
			\end{enumerate}

		\subsection{The Main Interface}
			Upon launching Dia, you will see a straightforward interface that separates your tools from your drawing canvas, as shown in \Cref{fig:DiaMain}.
			\begin{figure}[htbp]
				\centering
				\includegraphics[width=0.6\linewidth]{example-image}
				\caption{Dia Main Interface showing the Toolbox, Shape Selector, and Diagram Canvas.}
				\label{fig:DiaMain}
			\end{figure}
			
			The primary areas include:
			\begin{itemize}
				\item \textbf{Toolbox:} Located on the left, containing basic drawing tools, text insertion, and line connectors.
				\item \textbf{Shape Selector:} Located just below the toolbox, this drop-down menu allows you to switch between different libraries (\textit{e.g.}, from `Flowchart' to `Logic').
				\item \textbf{Canvas:} Your main drawing area, complete with a background grid.
			\end{itemize}

	\section{Dia in a \LaTeX\ Workflow}
		The most crucial aspect of using Dia is how you get your completed diagram out of the software and into your \LaTeX\ thesis without losing quality. 
		Because diagrams are made of sharp lines and text, they must be exported as vector graphics.

		\subsection{Exporting as Encapsulated PostScript (EPS)}
			For a very long time, EPS was the gold standard for \LaTeX\ vector graphics. 
			\begin{enumerate}
				\item In Dia, go to `File' \(\rightarrow\) `Export'.
				\item Select \textbf{Encapsulated PostScript (*.eps)} from the drop-down format list.
				\item This file can be compiled directly into your document. 
					However, note that modern \LaTeX\ workflows using \texttt{pdflatex} often convert EPS to PDF on the fly during compilation.
			\end{enumerate}

		\subsection{Exporting as PDF via Cairo}
			For modern uAlberta templates that rely on \texttt{pdflatex}, exporting directly to a PDF is the safest and cleanest method.
			\begin{enumerate}
				\item Go to `File' \(\rightarrow\) `Export'.
				\item Choose the \textbf{Cairo Portable Document Format (*.pdf)} option.
				\item This creates a tightly cropped, infinitely scalable PDF of your diagram that you can insert using the standard \cmd{includegraphics} or the custom \cmd{insertimage} command.
			\end{enumerate}

		\subsection{Advanced: Exporting to PGF/TikZ}
			Just like GNU Octave, Dia has a native export filter for \LaTeX's drawing package, \pkg{TikZ}. 
			If you want the text in your diagram to perfectly match the font of your thesis (\textit{e.g.}, Times New Roman or Computer Modern), you can export your diagram as a \textbf{PGF macro file (*.tex)}. 
			
			This generates raw \LaTeX\ code. 
			You can then include this diagram in your document using \cmd{insertplot}\mopt{my_diagram.tex}. 
			\note{%
				While exporting to \pkg{TikZ} ensures perfect font matching, complex Dia diagrams can generate a massive amount of raw code, which may slightly increase your document compile time. 
				For most students, exporting via the Cairo PDF method strikes the best balance between quality and ease of use.}